\section{Results and Discussion}\label{sec:results}

We report the primary endpoint, the mechanism-level and cost results, the
TMR-style targeting effect, the sim-vs-real agreement, the realised power, and
the bias controls. \textbf{Every number below is reproduced from the committed
Phase~5 result by a script; all of it is a mechanism demonstration under the
deterministic mock LLM} (model \texttt{\SWmodelID}, mocked=\SWmodelMocked;
git \texttt{\SWgitCommit}) --- not a claim about a real Claude model
(Section~\ref{sec:limitations}). Per-arm metrics for the primary
(\SWprimaryRegime) regime are in Table~\ref{tab:arms}.

\begin{table}[t]\centering
  \caption{Per-arm metrics in the primary (\SWprimaryRegime{}) regime: final
  average accuracy (ACC), backward and forward transfer (BWT, FWT),
  mechanism-level prune precision/recall/F1 against ground-truth relevance,
  final memory-vector count, and total tokens (thousands). Means over
  $n=\SWnSeeds$ seeds. The \texttt{oracle} arm (prune exactly the injected
  distractors) and \texttt{long\_context} (no pruning) are ceilings/baselines,
  not deployable treatments. \MockCaveat}
  \label{tab:arms}
  \small
  % AUTO-GENERATED by slow_wave.paper.numbers -- DO NOT EDIT BY HAND.
% Per-arm metrics, primary regime, mean over seeds. Regenerate with:
%   python -m slow_wave.paper.numbers
\begin{tabular}{l rrr rrr rr}
\toprule
Arm & ACC & BWT & FWT & Prec. & Rec. & F1 & Vec. & kTok \\
\midrule
no-sleep & 0.271 & -0.911 & +0.000 & 0.644 & 0.693 & 0.668 & 60 & 1.1 \\
replay-only & 0.279 & -0.901 & +0.000 & 0.643 & 0.692 & 0.667 & 60 & 1.1 \\
downscale-only & 0.271 & -0.911 & +0.000 & 0.644 & 0.693 & 0.668 & 60 & 1.1 \\
full-dream & 0.650 & -0.438 & +0.000 & 0.748 & 0.507 & 0.604 & 144 & 5.8 \\
reflection & 0.754 & -0.307 & +0.000 & 0.769 & 0.470 & 0.583 & 148 & 3.4 \\
random-pruning & 0.150 & -0.484 & +0.000 & 0.646 & 0.845 & 0.733 & 30 & 1.1 \\
oracle & 0.458 & -0.677 & +0.000 & 0.783 & 1.000 & 0.878 & 34 & 1.0 \\
long-context & 1.000 & +0.000 & +0.000 & 0.000 & 0.000 & 0.000 & 200 & 1.1 \\
A/A & 0.271 & -0.911 & +0.000 & 0.644 & 0.693 & 0.668 & 60 & 1.1 \\
\bottomrule
\end{tabular}

\end{table}

\subsection{Primary endpoint}
On the preregistered primary endpoint --- the paired (by-seed) difference in
final ACC between \texttt{full\_dream} and \texttt{no\_sleep} in the
\SWprimaryRegime{} regime --- the verdict is \textbf{\SWprimaryVerdict}. The mean
paired difference is $\Delta=\SWprimaryDiff$ (bootstrap 95\% CI
$[\SWprimaryCIlo,\SWprimaryCIhi]$, \SWprimaryCImethod{} method), the paired
\SWprimaryTest{} test gives $p=\SWprimaryP$, and the standardised paired effect
is Cohen's $d=\SWprimaryD$ (\SWprimaryEffectMag). The effect exceeds the A/A
noise floor of \SWnoiseFloor{} (exceeded=\SWexceedsNoiseFloor), so it clears the
bar a seed-only difference would produce. The same direction holds in all three
regimes: $\Delta=\SWdiffSignalRich$ (signal-rich, \SWverdictSignalRich),
$\SWdiffBalanced$ (balanced, \SWverdictBalanced), and $\SWdiffDistractorHeavy$
(distractor-heavy, \SWverdictDistractorHeavy). Within the mock-LLM regime, then,
the full dream cycle's operators behave as designed: at matched budget,
consolidation transfers signal into the semantic store so it survives episodic
eviction, lifting final accuracy over the no-sleep reference.

The per-task retention pattern (Figure~\ref{fig:retention}) shows the mechanism:
the no-sleep baseline retains only the most recent tasks (strongly negative
backward transfer in Table~\ref{tab:arms}), whereas replay-bearing arms retain
earlier-task signal. The replay\,$\times$\,downscale $2\times2$ ablation
(Figure~\ref{fig:ablation}) isolates the contribution of each mechanism, and the
registered secondary contrast confirms that the full cycle also beats
replay-only at matched budget ($\Delta=\SWdreamVsReplayDelta$, 95\% CI
$[\SWdreamVsReplayCIlo,\SWdreamVsReplayCIhi]$, Wilcoxon $p=\SWdreamVsReplayP$) ---
i.e.\ in this regime the effect is not merely re-exposure.

\begin{figure}[t]\centering
  \includegraphics[width=0.82\linewidth]{retention_curves.pdf}
  \caption{Per-task retention (final-row accuracy $R[T{-}1,j]$ after the full
  stream) per arm in the \SWprimaryRegime{} regime. $n=\SWnSeeds$ seeds; shaded
  bands are 95\% CIs (normal approximation, mean $\pm 1.96\,\mathrm{SE}$ over
  seeds). \MockCaveat}
  \label{fig:retention}
\end{figure}

\begin{figure}[t]\centering
  \includegraphics[width=0.9\linewidth]{ablation_table.pdf}
  \caption{Replay\,$\times$\,downscale $2\times2$ ablation
  (\texttt{no\_sleep}/\texttt{replay\_only}/\texttt{downscale\_only}/\texttt{full\_dream})
  as a mean-ACC heatmap, beside per-arm mean ACC. $n=\SWnSeeds$ seeds; error
  bars are 95\% CIs (normal approximation over seeds). \MockCaveat}
  \label{fig:ablation}
\end{figure}

\subsection{Consolidation quality, decoupled from accuracy}
Because the bench knows ground truth offline, it scores pruning directly
(Figure~\ref{fig:mechanism}). The \texttt{oracle} arm attains the highest prune
precision/recall (it prunes exactly the injected distractors), as it must ---
a sanity ceiling. Critically, prune quality and downstream accuracy can diverge:
in Table~\ref{tab:arms} the highest-accuracy deployable arms are not the
highest-precision pruners, which is exactly why the bench reports the two
\emph{decoupled} rather than collapsing them into a single score.

\begin{figure}[t]\centering
  \includegraphics[width=0.82\linewidth]{mechanism_pr.pdf}
  \caption{Mechanism-level prune quality (precision/recall/F1) per arm in the
  \SWprimaryRegime{} regime, decoupled from downstream accuracy. $n=\SWnSeeds$
  seeds; error bars are 95\% CIs (normal approximation over seeds). \MockCaveat}
  \label{fig:mechanism}
\end{figure}

\subsection{Cost, and the long-context comparison (a first-class negative)}
The accuracy-vs-compute frontier (Figure~\ref{fig:pareto}) and the stream-length
sweep (Figure~\ref{fig:crossover}) carry the most important caveat to the
positive primary result. The \texttt{long\_context} baseline --- stuffing the
entire stream into context with no consolidation --- attains the highest raw
accuracy at the highest token and memory cost. On the cost-adjusted metric
(\SWcrossoverMetric{}, \SWcrossoverTreatment{} vs \SWcrossoverBaseline{}) there
is \textbf{no crossover} in the swept range $L\in\{\SWcrossoverLengths\}$
(crossover found=\SWcrossoverFound): the long-context baseline keeps its
raw-accuracy lead across all tested lengths. This is precisely one of our
preregistered negative-result forms --- ``long-context wins below stream length
$L$'' --- and we report it as a first-class finding, not a footnote: within the
lengths we could afford to sweep, external consolidation does \emph{not} beat
stuffing the stream into a long context on raw accuracy, and whether a crossover
exists at larger $L$ is an open, explicitly unanswered question. The
contribution here is to \emph{map the cost regime}: consolidation's advantage in
this bench is a matched-budget, memory-footprint advantage (the consolidated
arms hold far fewer vectors than \texttt{long\_context} in Table~\ref{tab:arms}),
not a raw-accuracy win over an unbounded context.

\begin{figure}[t]\centering
  \includegraphics[width=0.82\linewidth]{cost_pareto.pdf}
  \caption{Accuracy vs compute (total tokens) per arm in the \SWprimaryRegime{}
  regime; the Pareto-non-dominated frontier (max ACC, min tokens) is
  highlighted. $n=\SWnSeeds$ seeds; error bars are 95\% CIs (normal
  approximation over seeds). \MockCaveat}
  \label{fig:pareto}
\end{figure}

\begin{figure}[t]\centering
  \includegraphics[width=0.82\linewidth]{long_context_crossover.pdf}
  \caption{Long-context crossover: cost-adjusted accuracy (\SWcrossoverMetric{},
  the preregistered metric) for \SWcrossoverTreatment{} vs \SWcrossoverBaseline{}
  across stream length $L$. No cost-adjusted crossover appears in the swept
  range. $n=\SWnSeeds$ seeds; error bars are 95\% CIs (normal approximation over
  seeds). \MockCaveat}
  \label{fig:crossover}
\end{figure}

\subsection{TMR-style targeting}
Benchmarking the replay-targeting analogue against the sleep TMR meta-analysis
\citep{hu2020tmr}, replay-bearing arms show a mean signal-retention lift of
$\SWtmrLift$ over non-replay arms, a bias-corrected Hedges'
$g=\SWtmrG$ (95\% CI $[\SWtmrCIlo,\SWtmrCIhi]$), which exceeds the benchmark
$g\approx\SWtmrBench$ (Figure~\ref{fig:tmr}). We stress the proxy nature of this
comparison: this is a prioritized-replay \emph{targeting analogue}, not a literal
cued-TMR protocol, so the comparison is to a benchmark effect size, not a claim
of biological equivalence.

\begin{figure}[t]\centering
  \includegraphics[width=0.75\linewidth]{tmr_targeting.pdf}
  \caption{TMR-style targeting: signal-retention lift for replay arms over
  no-replay arms (a replay-targeting analogue, \emph{not} a literal cued-TMR
  protocol), against the \citet{hu2020tmr} benchmark Hedges' $g=\SWtmrBench$.
  $n=\SWnSeeds$ seeds; the $g$ error bar is a bootstrap 95\% CI. \MockCaveat}
  \label{fig:tmr}
\end{figure}

\subsection{Sim-vs-real agreement, power, and bias controls}
Time-compression does not invert the arm ranking (Figure~\ref{fig:simreal}):
across key arms the accelerated high-compression sim and the real
low-compression long-horizon runs agree with Pearson $\SWsimPearson$ and
Spearman $\SWsimSpearman$ (\SWsimInversions{} inversions), with a maximum
absolute mean-ACC divergence of $\SWsimMaxDiv$. The ranking is preserved, so the
accelerated sim is a valid proxy for ordering arms here; the absolute divergence
is itself reported rather than hidden. For \emph{power}, the realised
$n=\SWpowerN$ seeds clears the floor of \SWpowerFloor; at the observed effect
($d=\SWpowerObsD$) a paired design requires only $n\ge\SWpowerReqN$ seeds
\citep{colas2018seeds}, so the design is amply powered \emph{for an effect this
large} --- a caveat we revisit in Section~\ref{sec:limitations}, since such a
large $d$ is itself an artifact of the deterministic mock regime. The bias
controls are clean: the temperature-0 stability control finds the mock
summariser identical across repeats (token CV $=\SWstabilityCV$,
identical=\SWstabilityIdentical), and the memory-drift detector reports
faithfulness $\SWdriftFaithfulness$ below the degradation threshold
$\SWdriftThreshold$ (degraded=\SWdriftDegraded).

\begin{figure}[t]\centering
  \includegraphics[width=0.82\linewidth]{sim_vs_real.pdf}
  \caption{Sim-vs-real agreement: per-arm mean ACC on the accelerated
  high-compression sim stream vs the real low-compression long-horizon stream,
  annotated with rank agreement and any inversion (Spearman $=\SWsimSpearman$;
  ranking preserved). $n=\SWnSeeds$ seeds; error bars are 95\% CIs (normal
  approximation over seeds). \MockCaveat}
  \label{fig:simreal}
\end{figure}

\paragraph{Interpretation.} In the mock-LLM mechanism-demonstration regime, the
apparatus does what it was built to do: each operator is independently
controllable, consolidation provably moves signal into a durable store, prune
quality is measured against ground truth and is decoupled from accuracy, and the
honest cost picture (long-context dominance on raw accuracy) is surfaced rather
than buried. What these results do \emph{not} establish --- and cannot, under a
mock LLM --- is that a real Claude model would show the same effect, or that the
effect would survive a real summariser's nondeterminism and drift. That is the
experiment the bench is built to run next.
